\scanpage{29}
\textbf{\mspage{214} -- Série I}\hfill\textbf{C. R. Acad. Sc. Paris, t. 292 (19 janvier 1981)}
\figure{\begin{tikzpicture}\draw (0,0)--(14,0);\end{tikzpicture}}

($b$) (\textit{Séparation}). $\Omega$ est un faisceau « (topologiquement) séparé », c'est-à-dire pour toute variété $\mathrm V$ ($\dim\mathrm V\leqq\mathrm N$) et tout couple $(\omega,\omega')$ d'orientations de $\mathrm V$ l'ensemble ouvert $\mathrm C_{\omega,\omega'}=\{x\in\mathrm V\mid\omega(x)=\omega'(x)\}$ est aussi fermé.

($c$) (\textit{Isomorphisme local}). Pour toute variété à bord $\mathrm V$ ($\dim\mathrm V\leqq\mathrm N$) l'application de Stokes $\delta_{\mathrm V}$ est un « isomorphisme d'espaces étalés sur $\mathrm V$ et $\partial\mathrm V$ », c'est-à-dire les conditions ci-dessous sont vérifiées :

(i) Si $\omega,\omega'\in\Omega^\bullet(\mathrm V)$ « induisent » la même orientation sur le bord $\partial\mathrm V$ de $\mathrm V$ [i. e. $\delta_{\mathrm V}(\omega)=\delta_{\mathrm V}(\omega')$], alors $\omega$ et $\omega'$ coïncident sur un voisinage ouvert de $\partial\mathrm V$.

(ii) Si $\bar\omega\in\Omega(\partial\mathrm V)$, alors $\bar\omega$ est « localement induite », c'est-à-dire pour tout $x\in\partial\mathrm V$ il existe un voisinage ouvert $\mathrm U$ de $x$ (en $\mathrm V$) et une orientation $\omega\in\Omega^\bullet(\mathrm U)$ tel que $\delta_{\mathrm V}(\omega)=\bar\omega\mid\partial\mathrm U$ ($\bar\omega\mid\partial\mathrm U$ désigne la restriction de $\bar\omega$ à $\partial\mathrm U$).

On voit aussitôt que pour toute variété à bord $\mathrm V$ la conjonction des conditions (i) et (ii) ci-dessus équivaut à dire que l'application :
\diagram{\begin{tikzcd}\Omega^\bullet_{\mathrm V}\mid\partial\mathrm V\arrow[r,"\sim"]&\Omega_{\partial\mathrm V},\end{tikzcd}}
déduite de l'application de Stokes est un \textit{isomorphisme d'espaces étalés de base} $\partial\mathrm V$.

Avec les données et axiomes précédents on démontre la proposition suivante :

\textsc{Proposition}. -- \textit{Pour tout objet $\mathrm V$ de $\mathcal V^\bullet_{\leqq\mathrm N}$ l'espace étalé $\Omega_{\mathrm V}$ sur $\mathrm V$ défini par $\Omega$ est un revêtement à deux feuillets (appelé} revêtement des orientations de $\mathrm V$).

Il résulte de cette proposition que le groupe $\mu_2=\{\pm1\}$ opère sur le faisceau $\Omega$. Si $\omega$ est une orientation d'une variété $\mathrm V$ on pose $-\omega=-1.\omega$; l'orientation $-\omega$ est appelée l'\textit{orientation opposée à} $\omega$.

($d$) (\textit{Anti-commutativité}). -- Pour tout $\mathrm V$ objet de $\mathcal V_{\leqq\mathrm N}$ et pour des sous-variétés $\mathrm V'$, $\mathrm V''$ de $\mathrm V$ tels que $\mathrm V'\cup\mathrm V''=\mathrm V$ et $\mathrm V'\cap\mathrm V''=\partial\mathrm V'=\partial\mathrm V''$, le diagramme :
\diagram{\begin{tikzcd}[column sep=large,row sep=small]
&\Omega^\bullet(\mathrm V')\arrow[dr,"\delta_{\mathrm V'}"]&\\
\Omega(\mathrm V)\arrow[ur,"r'"]\arrow[dr,"r''"']&&\Omega(\partial\mathrm V')=\Omega(\partial\mathrm V''),\\
&\Omega^\bullet(\mathrm V'')\arrow[ur,"\delta_{\mathrm V''}"']&
\end{tikzcd}}
\textit{où} $r'$ et $r''$ sont les restrictions, est « anti-commutatif », c'est-à-dire pour toute orientation $\omega$ de $\mathrm V$ on a la relation :
\[
\delta_{\mathrm V'}(\omega\mid\mathrm{Int}(\mathrm V'))=-\delta_{\mathrm V''}(\omega\mid\mathrm{Int}(\mathrm V'')).
\]

\textsc{Définition}. -- Soient $\mathrm T=(\Omega,\delta,\nu)$ et $\mathrm T'=(\Omega',\delta',\nu')$ deux théories de l'orientation des variétés topologiques en dimension $\leqq\mathrm N$. Un morphisme de $\mathrm T$ en $\mathrm T'$ est un morphisme fonctoriel de $\Omega$ en $\Omega'$ qui commute aux données de Stokes et de normalisation, en un sens évident.

\textsc{Théorème} 1 (\textit{cf.} [3]). -- ($a$) \textit{Il existe une théorie de l'orientation des variétés topologiques en dimension $\leqq N$.}

($b$) \textit{Il existe au plus \textup{un morphisme} entre deux théories de l'orientation des variétés topologiques en dimension $\leqq\mathrm N$ et un tel morphisme est un \textup{isomorphisme}.}

($c$) \textit{Si $\mathrm N\ne4$ il existe un \textup{(unique)} isomorphisme entre deux théories de l'orientation des variétés topologiques en dimension $\leqq\mathrm N$.}

On démontre ce théorème en utilisant les résultats ci-après.
\margin{28}
